\section{Introduction}
\label{sec:introduction}

\subsection{Background and motivation}

揚水発電は, 変動性再生可能エネルギーの導入拡大に伴い, 電力系統の需給調整と短周期のエネルギー貯蔵を担う基盤技術として, その重要性を一層高めている \cite{PSP_background1,PSP_background2}. とくに, 大規模電力系統においては, 出力変動の大きい風力・太陽光発電を補完し, 短時間での出力追従と系統安定化を可能にする手段として,  \cite{PSP_background3,PSP_background4}.

一方で, 揚水発電の上部調整池(以下, 上池)は, 取放水を繰り返す運用に長期間曝されるため, 池内流況とそれに起因する堆砂・濁質滞留が, 運用・維持管理上の重要課題となる \cite{PSP_sed1,PSP_sed2,PSP_sed3}. とくに, 上池が浅く広い形状を有する場合には, 自由水面をもつ浅水流としての性質が強くなり, 平面内にはジェットの偏向, 再循環, 再付着, 低流速域の形成といった特徴的な流動構造が現れうる. これらの構造は, 堆積・滞留の空間分布だけでなく, 取放水構造物近傍の流況変動, さらには数値解析時の境界条件設定や検証方針にも影響する \cite{shallow_reservoir1,shallow_reservoir2}.

このような背景のもと, 上池内流れを実務的に理解し, 将来的な三次元数値解析(RANS/LES など)へ接続するためには, まず支配的な平面内流動構造を, 再現性のある条件設定の下で定量化し, 流動レジームとして整理することが重要である. しかし, 実機上池は大規模であり, 現地における詳細な流況観測は容易ではない. そのため, 実機条件を念頭に置いた室内模型実験により, 上池内流れの骨格を把握し, 後続の数値解析や設計検討に資する基礎データを整備することに意義がある.

\subsection{Related work and gap}

浅い矩形貯水池や矩形タンクにおける流れについては, sudden expansion や浅い再循環流の文脈で, 入口ジェットの偏向, 側壁への再付着, 対称・非対称な循環構造, およびそれらの遷移挙動が実験的に調べられてきた \cite{shallow_recirc1,shallow_recirc2,shallow_recirc3}. とくに Dufresne ら \cite{Dufresne2010} は, 浅い矩形貯水池における流動パターンを体系的に分類し, 対称流, 非対称流, および不安定な遷移流を整理するとともに, 幾何条件と水理条件に基づくレジーム遷移の枠組みを提示している. これらの知見は, 浅い上池模型における表面流構造を整理する上で重要な基礎を与える.

しかし, 揚水発電上池の取放水は, 一般の一方向流入・一方向流出系とは異なり, 同一の取放水構造物が inflow と outflow の両方を担う可逆系である \cite{PSP_io1,PSP_io2}. この可逆性は, basin scale の流れに対して少なくとも二つの含意をもつ. 第一に, inflow と outflow は幾何学的には同一開口を共有していても, 流れ場としては必ずしも対称にならず, basin 内の循環構造や低流速域の分布が大きく異なりうる. 第二に, outflow 側では平均場だけでは表現できない非定常性が強く現れる可能性がある. 実際, Zhu ら \cite{Zhu2023} は, 揚水発電所の inlet/outlet 構造物近傍において, outflow 条件では瞬時流速の変動が大きく, 時間平均流速だけでは流況を十分に表現できない一方, inflow 条件ではその変動が比較的小さいことを示している. このことは, 可逆取放水系の評価において, 平均ベクトル場のみでは不十分であり, time-resolved な観測量が必要であることを示唆する.

% NOTE (2026-07, 解析結果との整合): 下段落の予備観察「outflow では basin 全体を
% 占める時計回り循環」は, 静水開始の本計測 (24 runs) では確認されなかった.
% 静水開始では outflow は回転をもたない対称収束シンクであり, 予備観察の循環は
% 先行 inflow の残留 (履歴効果, Muller 2012) と解釈される. Results/Discussion
% はその立場で書いたので, 本段落の予備観察の記述もそれに合わせて
% 「連続運転時の予備観察では〜」等に修正することを推奨.
さらに, 上池全体を対象とした basin-scale の模型実験では, 取放水構造物そのものの局所流況ではなく, basin 内の大域的な循環構造, 低流速域, およびその遷移を捉える必要がある. このとき, 水深と開口高さの相対関係は, 表面流の支配様式を決める重要な因子になりうる. 実際, 本研究に先行する予備的な模型実験では, inflow では偏向したジェットと二つの循環構造, outflow では basin 全体を占める時計回り循環が観察され, さらに水深が開口高さ付近に達したときに流況の様相が変化する兆候が確認されている \cite{Dufresne2010,Zhu2023}. 一方で, 可逆単一開口を有する浅い basin を対象として, inflow/outflow を独立した初期条件のもとで比較し, 相対水深 $h/a$ によるレジーム遷移と outflow の非定常性を表面計測に基づいて整理した研究は, 著者らの知る限り十分ではない.

したがって, 本研究のギャップは, 可逆単一開口を有する浅い矩形 basin において, basin-scale の表面流レジームを, inflow/outflow の非対称性, 相対水深 $h/a$, および time-resolved な非定常指標の観点から体系的に整理した知見が不足している点にある.

\subsection{Objectives and contributions}

本研究は, 揚水発電上池を想定した浅い矩形水槽において, 単一の可逆開口を介した inflow および outflow の表面流動を, time-resolved な表面PIVにより計測し, 水深および流量条件に対する流動レジーム遷移を定量的に整理することを目的とする. ここで本稿のスコープは flow-only に限定し, 堆砂そのものの再現や内部流速分布の直接計測は扱わない. その代わり, 将来的な堆砂実験および三次元数値解析に接続するための, basin-scale の表面流動構造の整理を第一段階として位置付ける.

本研究で対象とする主要な問いは, 次の3点である.
(1) 同一開口を共有する inflow と outflow は, basin 内表面流のレジームとしてどの程度非対称であるか.
(2) 開口高さ $a$ に対する相対水深 $h/a$ は, 表面流レジームの遷移をどのように支配するか.
(3) とくに outflow において, 平均場では捉えにくい非定常性は, どのような time-resolved 指標によって特徴付けられるか.

以上を踏まえた本研究の貢献は, 次のように要約できる.
\begin{itemize}
  \item 単一可逆開口を有する浅い矩形水槽に対して, inflow と outflow を独立条件で計測し, basin-scale の表面流レジームを水深および流量の関数として体系的に整理する.
  \item 相対水深 $h/a$ に着目し, レジーム遷移が生じやすい条件帯を特定するとともに, inflow/outflow 間の流動非対称性を明確化する.
  \item time-resolved な表面PIV解析に基づき, 低流速域面積率, 循環指標, 速度変動指標などのスカラー量を導入し, outflow の非定常性を平均場を超えて記述する.
  \item 実機上池を念頭に置いた室内模型実験として, 将来的な堆砂研究および三次元数値解析の基礎となる表面流データセットとレジーム整理の枠組みを提示する.
\end{itemize}
